% =============================================================
% CAPÍTULO 3 — Metodología
% Versión limpia: sin negritas en cuerpo de texto, con figura
% de pipeline y expansiones integradas.
%
% Requisitos en el preámbulo (añadir si no están):
%   \usepackage{tikz}
%   \usetikzlibrary{arrows.meta, positioning, fit, backgrounds, calc}
% =============================================================

\section{Metodología}
\label{sec:metodologia-tfg}

Este capítulo describe el diseño experimental del proyecto: el dataset y su
preprocesamiento (\S\ref{sec:met-dataset}), el modelo base y su \emph{fine-tuning}
(\S\ref{sec:met-modelo}), el módulo de compresión SVD (\S\ref{sec:met-svd}),
las cuatro técnicas de interpretabilidad mecánica empleadas
(\S\ref{sec:met-interpretabilidad}) y las métricas de evaluación
(\S\ref{sec:met-metricas}). La Figura~\ref{fig:pipeline} muestra el flujo
experimental completo.

% ── Figura de pipeline ──────────────────────────────────────

\begin{figure}[H]
    \centering
    \begin{tikzpicture}[
        box/.style={
            rectangle, rounded corners=3pt, draw=black!35,
            minimum height=0.9cm, minimum width=2.6cm,
            align=center, font=\small,
            inner xsep=8pt, inner ysep=5pt
        },
        smallbox/.style={
            rectangle, rounded corners=2pt, draw=black!25,
            minimum height=0.7cm, align=center, font=\scriptsize,
            inner sep=4pt, fill=gray!6
        },
        arr/.style={-{Stealth[length=4pt]}, semithick, black!45},
        dasharr/.style={-{Stealth[length=4pt]}, thin, black!30, dashed},
        lbl/.style={font=\scriptsize\itshape, text=black!40},
    ]

    % Row 1 (y=0): GoEmotions → BERT → Baseline
    \node[box, fill=gray!8]  (data)     at (0,   0) {GoEmotions\\[-1pt]{\scriptsize 23 emociones}};
    \node[box, fill=blue!6]  (model)    at (3.6, 0) {BERT-base\\[-1pt]{\scriptsize Fine-tuning}};
    \node[box, fill=teal!8]  (baseline) at (7.2, 0) {Baseline\\[-1pt]{\scriptsize F1\textsubscript{macro} = 0{,}577}};

    \draw[arr] (data) -- (model);
    \draw[arr] (model) -- (baseline);

    % Row 2 (y=-1.6): Compresión SVD
    \node[box, fill=orange!8] (svd) at (3.6, -1.6) {Compresión SVD\\[-1pt]{\scriptsize Uniforme / Adaptativa}};
    \draw[arr] (baseline.south) -- ++(0,-0.35) -| (svd.north);

    % Row 3 (y=-3.5): Interpretabilidad mecánica — manual rectangle
    \draw[rounded corners=4pt, draw=black!15, fill=blue!2]
        (-1.6, -2.85) rectangle (9.0, -4.55);
    \node[font=\small, text=black!60] at (3.6, -3.15) {Interpretabilidad mecánica};

    \node[smallbox] (probing)  at (0,   -3.85) {Probing\\[-1pt]{\tiny NB04 · capa}};
    \node[smallbox] (patching) at (2.6, -3.85) {Act.\ patching\\[-1pt]{\tiny NB05 · componente}};
    \node[smallbox] (ablation) at (5.2, -3.85) {Head ablation\\[-1pt]{\tiny NB06 · cabeza}};
    \node[smallbox] (neurons)  at (7.8, -3.85) {Neuronas FFN\\[-1pt]{\tiny NB07 · neurona}};

    \draw[arr] (svd.south) -- ++(0, -0.7);

    % Row 4 (y=-5.6): Compresión informada
    \node[box, fill=teal!10] (informed) at (3.6, -5.6) {Compresión informada\\[-1pt]{\scriptsize NB09 · Greedy data-driven}};
    \draw[arr] (3.6, -4.55) -- (informed.north);

    % Row 5 (y=-7.1): Fine-tuning recovery
    \node[box, fill=green!8] (recovery) at (3.6, -7.1) {Fine-tuning recovery\\[-1pt]{\scriptsize F1 = 0{,}591 (+2{,}4\%)}};
    \draw[arr] (informed) -- (recovery);

    % Feedback loop (right side)
    \coordinate (fb1) at (9.6, -5.6);
    \coordinate (fb2) at (9.6, -1.6);
    \draw[dasharr, rounded corners=3pt]
        (informed.east) -- (fb1) -- (fb2) -- (svd.east);
    \node[lbl, rotate=90, anchor=south] at (9.6, -3.6) {retroalimentación};

    \end{tikzpicture}
    \caption{Pipeline experimental. El modelo baseline se somete a compresión SVD y a
    cuatro técnicas de interpretabilidad mecánica con granularidad creciente (de capa
    completa a neurona individual). Los hallazgos retroalimentan el diseño de
    estrategias de compresión informadas (línea discontinua). El ciclo se cierra
    con \emph{fine-tuning} del modelo comprimido.}
    \label{fig:pipeline}
\end{figure}

% -------------------------------------------------------------
\subsection{Dataset: GoEmotions}
\label{sec:met-dataset}
% -------------------------------------------------------------

El dataset GoEmotions \cite{demszky2020goemotions} contiene 58.009 comentarios de
Reddit anotados por tres evaluadores con una taxonomía de 28 categorías emocionales
(27 emociones + \textit{neutral}). A diferencia de la clasificación de sentimientos
estándar (positivo/negativo), GoEmotions es un dataset de clasificación
multi-etiqueta: cada texto puede expresar múltiples emociones simultáneamente (por
ejemplo, \textit{admiración} y \textit{alegría}), lo que impone al modelo la necesidad
de resolver co-ocurrencias con matices semánticos cercanos.

\subsubsection{Filtrado de emociones}

Del conjunto original de 28 categorías, se eliminaron cinco que presentaban problemas
de rendimiento o representatividad:

\begin{itemize}[nosep]
    \item \textit{neutral}: excluida por no ser una emoción propiamente dicha y por
    dominar el dataset (14.000+ ejemplos), lo que distorsionaba las métricas
    macro-promediadas.
    \item \textit{grief}, \textit{nervousness}, \textit{pride}, \textit{relief}:
    excluidas por producir F1 = 0 en el baseline, debido a su bajísima frecuencia en
    el dataset (menos de 100 ejemplos cada una en el conjunto de entrenamiento). La
    presencia de emociones con F1 nulo degradaba el F1 macro sin aportar información
    útil sobre la capacidad del modelo.
\end{itemize}

El filtrado se implementó a nivel del pipeline de datos: se eliminaron las cinco
columnas correspondientes del vector de etiquetas (pasando de 28 a 23 dimensiones) y
se descartaron los ejemplos cuyo vector resultante quedaba vacío (aquellos anotados
exclusivamente con alguna de las emociones excluidas). El resultado es un dataset de
23 emociones donde todas las categorías producen un F1 funcional (mínimo: 0{,}267
para \textit{embarrassment}).

La Tabla~\ref{tab:emociones-23} lista las 23 emociones del dataset final, ordenadas
por F1 en el baseline.

\begin{table}[H]
    \centering
    \caption{Las 23 emociones del dataset final, con F1 en el baseline.}
    \label{tab:emociones-23}
    \small
    \begin{tabular}{@{}l r r l@{}}
        \toprule
        Emoción & F1 baseline & Ejs.\ train & Grupo semántico \\
        \midrule
        gratitude      & 0{,}927 & 2\,662 & Positiva social \\
        amusement      & 0{,}853 & 2\,328 & Positiva energética \\
        love           & 0{,}816 & 2\,086 & Positiva social \\
        admiration     & 0{,}757 & 4\,130 & Positiva social \\
        fear           & 0{,}710 & \phantom{0\,}596 & Negativa interna \\
        curiosity      & 0{,}707 & 2\,191 & Epistémica \\
        remorse        & 0{,}672 & \phantom{0\,}545 & Negativa interna \\
        joy            & 0{,}635 & 1\,452 & Positiva energética \\
        sadness        & 0{,}632 & 1\,326 & Negativa interna \\
        optimism       & 0{,}573 & 1\,581 & Positiva orientada \\
        surprise       & 0{,}571 & 1\,060 & Epistémica \\
        approval       & 0{,}555 & 2\,939 & Positiva social \\
        anger          & 0{,}550 & 1\,567 & Negativa reactiva \\
        caring         & 0{,}542 & 1\,087 & Positiva orientada \\
        desire         & 0{,}535 & \phantom{0\,}641 & Positiva orientada \\
        disapproval    & 0{,}527 & 2\,022 & Negativa reactiva \\
        disgust        & 0{,}521 & \phantom{0\,}793 & Negativa reactiva \\
        confusion      & 0{,}515 & 1\,368 & Epistémica \\
        excitement     & 0{,}454 & \phantom{0\,}853 & Positiva energética \\
        annoyance      & 0{,}411 & 2\,470 & Negativa reactiva \\
        disappointment & 0{,}275 & 1\,269 & Negativa interna \\
        realization    & 0{,}271 & 1\,110 & Epistémica \\
        embarrassment  & 0{,}267 & \phantom{0\,}303 & Negativa interna \\
        \midrule
        \multicolumn{2}{@{}l}{F1 macro promedio} & & 0{,}577 \\
        \bottomrule
    \end{tabular}
    \\[0.15cm]{\small Grupos semánticos propuestos a partir del clustering neuronal
    (Capítulo~\ref{sec:resultados-interpretabilidad}). Fuente: elaboración propia
    sobre GoEmotions \cite{demszky2020goemotions}.}
\end{table}

\subsubsection{Splits y preprocesamiento}

Se utilizaron los splits oficiales de GoEmotions: entrenamiento (43.410 ejemplos),
validación (5.426) y test (5.427), tras el filtrado. El preprocesamiento se limita a
la tokenización con el tokenizer de BERT-base-uncased, con una longitud máxima de
128 tokens y padding dinámico. No se aplican técnicas de aumento de datos.

% -------------------------------------------------------------
\subsection{Modelo: BERT-base}
\label{sec:met-modelo}
% -------------------------------------------------------------

\subsubsection{Arquitectura}

BERT-base-uncased \cite{devlin2019bert} consta de 12 capas de encoder Transformer,
cada una con un bloque de atención multi-cabeza (12 cabezas) y un bloque
feed-forward. El modelo tiene 110 millones de parámetros, de los cuales
84{,}9 millones corresponden a las 72 capas lineales del encoder (6 matrices por
capa $\times$ 12 capas), que constituyen el objetivo de la compresión SVD. La
Tabla~\ref{tab:capas-bert} detalla las dimensiones de cada tipo de componente.

\begin{table}[H]
    \centering
    \caption{Capas lineales del encoder de BERT-base (por capa).}
    \label{tab:capas-bert}
    \small
    \begin{tabular}{@{}l l c r@{}}
        \toprule
        Bloque & Componente & Dimensiones & Parámetros \\
        \midrule
        Atención & Query ($W_Q$)            & $768 \times 768$   & 589.824 \\
        Atención & Key ($W_K$)              & $768 \times 768$   & 589.824 \\
        Atención & Value ($W_V$)            & $768 \times 768$   & 589.824 \\
        Atención & Output ($W_O$)           & $768 \times 768$   & 589.824 \\
        FFN      & Intermediate ($W_{int}$) & $768 \times 3072$  & 2.359.296 \\
        FFN      & Output ($W_{out}$)       & $3072 \times 768$  & 2.359.296 \\
        \midrule
        \multicolumn{3}{l}{Total por capa del encoder} & 7.077.888 \\
        \multicolumn{3}{l}{Total 12 capas} & 84.934.656 \\
        \bottomrule
    \end{tabular}
    \\[0.15cm]{\small Fuente: elaboración propia.}
\end{table}

\subsubsection{\emph{Fine-tuning}}

El modelo se fine-tuneó sobre el dataset GoEmotions (23 emociones) añadiendo una capa
de clasificación lineal sobre el token \texttt{[CLS]} de la última capa del encoder.
La clasificación multi-etiqueta se implementa con una función de pérdida
\texttt{BCEWithLogitsLoss} (binary cross-entropy con logits), que trata cada emoción
como una predicción binaria independiente.

Los hiperparámetros de \emph{fine-tuning} se resumen en la Tabla~\ref{tab:hyperparams}.

\begin{table}[H]
    \centering
    \caption{Hiperparámetros de \emph{fine-tuning}.}
    \label{tab:hyperparams}
    \small
    \begin{tabular}{@{}l l@{}}
        \toprule
        Parámetro & Valor \\
        \midrule
        Optimizador        & AdamW \\
        Learning rate      & $2 \times 10^{-5}$ \\
        Epochs             & 4 \\
        Batch size         & 32 \\
        Warmup             & 10\% de los steps \\
        Weight decay       & 0{,}01 \\
        Max sequence length & 128 tokens \\
        Semilla aleatoria  & 42 \\
        \bottomrule
    \end{tabular}
    \\[0.15cm]{\small Fuente: elaboración propia.}
\end{table}

El modelo resultante alcanza un F1 macro de 0{,}577 y un F1 micro de 0{,}627 sobre
el conjunto de test, lo que constituye el baseline contra el que se mide toda
degradación posterior.

% -------------------------------------------------------------
\subsection{Compresión SVD}
\label{sec:met-svd}
% -------------------------------------------------------------

\subsubsection{Fundamento teórico}

La SVD factoriza cualquier matriz $W \in \mathbb{R}^{m \times n}$ como
$W = U \Sigma V^T$, donde $U \in \mathbb{R}^{m \times m}$ y
$V \in \mathbb{R}^{n \times n}$ son matrices ortogonales y
$\Sigma \in \mathbb{R}^{m \times n}$ es diagonal con los valores singulares
$\sigma_1 \geq \sigma_2 \geq \cdots \geq \sigma_{\min(m,n)} \geq 0$. El
teorema de Eckart-Young \cite{eckart1936approximation} garantiza que la mejor
aproximación de rango $k$ en norma de Frobenius es la SVD truncada:
\begin{equation}
    W_k = \sum_{i=1}^{k} \sigma_i u_i v_i^T = U_k \Sigma_k V_k^T
    \label{eq:svd-truncada}
\end{equation}
con error $\|W - W_k\|_F^2 = \sum_{i=k+1}^{\min(m,n)} \sigma_i^2$. La SVD induce
una descomposición ortogonal del espacio de salida $\mathbb{R}^m$:
\begin{equation}
    \mathbb{R}^m = \mathcal{S}_k \oplus \mathcal{S}_k^\perp,
    \quad \mathcal{S}_k = \text{span}\{u_1, \ldots, u_k\}
    \label{eq:subespacios}
\end{equation}
donde $\mathcal{S}_k$ es el \emph{subespacio conservado} (asociado a los $k$
mayores valores singulares) y $\mathcal{S}_k^\perp$ el \emph{subespacio
eliminado}. El operador de proyección sobre $\mathcal{S}_k^\perp$ es
$P_k^\perp = I - U_k U_k^T$, y la fracción de señal $x \in \mathbb{R}^m$
eliminada por la compresión a rango $k$ es
$\|P_k^\perp x\|_2 / \|x\|_2$. Esta cantidad se utiliza en
\S\ref{sec:neuron-analysis} para evaluar si la SVD discrimina selectivamente
entre direcciones funcionalmente significativas (vectores de selectividad
neuronal) o si elimina señal uniformemente.

En la práctica, una capa lineal $W$ de dimensión $m \times n$ se reemplaza
por dos capas: $V_k^T$ ($k \times n$) y $U_k \Sigma_k$ ($m \times k$),
reduciendo los parámetros de $mn$ a $k(m+n)$.

\subsubsection{Rango de break-even}

La compresión SVD solo reduce parámetros si el rango elegido es menor que
$mn/(m+n)$. Para BERT-base:
\begin{itemize}[nosep]
    \item Componentes de atención ($768 \times 768$): break-even en rango 384.
    \item Componentes FFN ($768 \times 3072$): break-even en rango 614.
\end{itemize}
Rangos superiores a estos umbrales \emph{aumentan} el número total de parámetros.
Este detalle técnico es relevante para interpretar correctamente las estrategias
adaptativas de alta energía (Capítulo~\ref{sec:resultados-compresion}).

\subsubsection{Implementación}

El módulo de compresión (\texttt{src/compression/svd.py}) implementa una clase
\texttt{SVDLinear} que almacena las dos matrices factorizadas y realiza la
multiplicación en dos pasos durante el forward pass. La función
\texttt{apply\_svd\_compression()} acepta un modelo BERT, una especificación de
rangos por capa y componente, y devuelve el modelo comprimido. La SVD se computa
siempre en \texttt{float32} para evitar inestabilidades numéricas, independientemente
de la precisión del modelo.

\subsubsection{Estrategias de compresión evaluadas}

Se evaluaron cuatro familias de estrategias:

\begin{enumerate}[nosep]
    \item Uniforme: mismo rango $k \in \{32, 64, 128, 256, 384, 512\}$ para
    las 72 matrices del encoder.
    \item Adaptativa (energía): rango por matriz determinado por un umbral
    de energía espectral ($E(k) \geq \{80\%, 90\%, 95\%, 99\%\}$), donde
    $E(k) = \sum_{i=1}^k \sigma_i^2 / \sum_{i=1}^n \sigma_i^2$.
    \item Informada heurística: rangos asignados manualmente a partir de
    reglas derivadas de los resultados de interpretabilidad (proteger capas tardías,
    comprimir agresivamente Q/K). Tres niveles de agresividad (light, moderate,
    aggressive).
    \item Greedy (data-driven): algoritmo que construye un menú de movimientos
    de compresión (cada uno del tipo ``comprimir componente $X$ en capas $Y$ al
    rango $Z$'') y los ordena por eficiencia, definida como parámetros ahorrados
    por punto de F1 perdido. El algoritmo selecciona movimientos del más eficiente
    al menos eficiente hasta alcanzar el ratio de parámetros objetivo. El coste
    en F1 de cada movimiento proviene de datos empíricos de sensibilidad por
    componente y profundidad, obtenidos en los notebooks de compresión y
    \emph{activation patching}.
\end{enumerate}

% -------------------------------------------------------------
\subsection{Técnicas de interpretabilidad mecánica}
\label{sec:met-interpretabilidad}
% -------------------------------------------------------------

Se emplearon cuatro técnicas complementarias que cubren escalas progresivas de
granularidad: capa completa (probing), capa y componente (\emph{activation patching}),
cabeza individual (ablación) y neurona individual (análisis de selectividad).

\subsubsection{Probing lineal}
\label{sec:met-probing}

Para cada una de las 13 capas de BERT (embedding + 12 encoder layers) y cada
una de las 23 emociones, se entrena un probe de regresión logística
(\texttt{LogisticRegression} de scikit-learn, $C = 1{,}0$ sin tuning) sobre
la representación del token \texttt{[CLS]}. La decisión de fijar $C = 1{,}0$
sigue la lógica de Hewitt y Liang \cite{hewitt2019designing}: el probe debe
ser deliberadamente simple para medir la calidad de la representación, no la
capacidad del clasificador. Esta elección tiene una contrapartida que conviene
reconocer: información codificada no linealmente en capas tempranas no sería
detectada por un probe lineal, lo que podría infraestimar la presencia de
señal emocional en las primeras capas. La interpretación de los resultados
(\S\ref{sec:probing}) se realiza con esta cautela.

El total nominal de probes es $23 \times 13 = 299$; sin embargo, los probes
sobre la capa de embedding producen F1 = 0 uniformemente para todas las
emociones (el embedding crudo no contiene información emocional separable
linealmente), por lo que se descartan del análisis efectivo. El número
efectivo de probes que entran en las correlaciones y métricas es de
$23 \times 12 = 276$, valor que utilizamos consistentemente en el resto de la
memoria.

La representación utilizada es el vector del token \texttt{[CLS]} de cada
capa, en lugar del mean-pooling sobre todos los tokens. Esta decisión sigue
la convención del propio BERT: durante el \emph{fine-tuning}, el token \texttt{[CLS]}
es el único cuya representación se conecta directamente a la cabeza de
clasificación, por lo que es la posición donde el modelo concentra la
información relevante para la tarea. Utilizar mean-pooling mediría la
información distribuida en toda la secuencia, lo que reflejaría una capacidad
representacional distinta de la que el modelo realmente explota para
clasificar.

La métrica de probing es el F1 por probe. Se define la \emph{capa de
cristalización} de una emoción como la primera capa donde su F1 alcanza el
80\% del máximo observado en cualquier capa. Este umbral es una convención;
el orden de cristalización es estable con umbrales de 0{,}7 y 0{,}9.

\subsubsection{\emph{Logit lens}}
\label{sec:met-logitlens}

El \emph{logit lens} \cite{nostalgebraist2020logit} aplica la cabeza de
clasificación entrenada del modelo (\texttt{BertPooler} + clasificador
lineal) a la representación del token \texttt{[CLS]} de cada capa
intermedia, sin re-entrenamiento. Para cada capa $\ell \in \{0, \ldots, 11\}$
y cada ejemplo, se calcula:
\begin{equation}
    \hat{p}_\ell = \sigma\left(W_{\text{cls}} \cdot \tanh\left(W_{\text{pool}} \cdot h_\ell^{[\text{CLS}]}\right) + b_{\text{cls}}\right) \in \mathbb{R}^{23}
    \label{eq:logitlens}
\end{equation}
donde $h_\ell^{[\text{CLS}]}$ es la representación del token \texttt{[CLS]}
en la capa $\ell$ y $\sigma$ es la sigmoide aplicada componente a
componente. Las predicciones se agregan sobre 2\,300 ejemplos del conjunto
de test (submuestreo equilibrado entre los tres grupos de cristalización del
probing). Las métricas reportadas son la sigmoid media de la emoción
correcta (\emph{gold sigmoid}), la sigmoid media de la emoción más probable
(\emph{top-1 sigmoid}) y la suma total de las 23 sigmoids.

Esta técnica tiene una limitación de calibración conocida que se discute en
\S\ref{sec:logit-lens}: los pesos del pooler y del clasificador fueron
optimizados exclusivamente para representaciones de la última capa tras
LayerNorm, por lo que los valores absolutos en capas intermedias no son
estrictamente comparables al baseline. Las interpretaciones se realizan
sobre tendencias relativas (entre capas o entre grupos emocionales), no
sobre magnitudes absolutas.

\subsubsection{\emph{Activation patching}}
\label{sec:met-patching}

El \emph{activation patching} parte de un modelo completamente colapsado:
comprimido con SVD uniforme a rango 64 (F1 = 0{,}000). Se eligió el rango 64
como punto de colapso porque es el rango más alto que produce un colapso
total y uniforme del modelo (F1 = 0{,}000 en las 23 emociones), garantizando
un estado base completamente destruido desde el cual cualquier recuperación
es atribuible al componente restaurado. Rangos inferiores (32) producen el
mismo colapso, pero con mayor distorsión de los pesos que podría introducir
artefactos en el patching. Rangos superiores (128) permiten recuperación
parcial espontánea en algunas emociones, lo que contaminaría la atribución
causal.

La corrupción mediante SVD truncada a rango 64 no es
\emph{neutral} en sentido estadístico (a diferencia del \emph{causal tracing}
clásico de Meng et al.\ \cite{meng2022locating}, que usa ruido gaussiano
sobre los embeddings), sino una distorsión \emph{estructurada} que respeta
la geometría espectral de los pesos. Las conclusiones extraídas tienen,
por tanto, naturaleza \emph{funcional} (qué componentes son
suficientes para reactivar el modelo desde el colapso) más que estrictamente
causal en el sentido de Pearl. Esta distinción se discute con detalle en
\S\ref{sec:activation-patching}; reconocerla forma parte del rigor
metodológico del trabajo.

Desde este estado, se restauran selectivamente los pesos originales de una
capa o componente y se mide la recuperación del F1 macro. La métrica de
interés es el porcentaje de restauración:
\begin{equation}
    \text{Restauración}(\%) = \frac{F1_{\text{patched}} - F1_{\text{comprimido}}}{F1_{\text{baseline}} - F1_{\text{comprimido}}} \times 100
    \label{eq:restauracion}
\end{equation}

El patching se realiza a dos niveles de granularidad: (a) capa completa,
restaurando las 6 matrices de una capa del encoder; y (b) componente,
restaurando solo el bloque de atención o solo el bloque FFN de una capa.

\subsubsection{Ablación de cabezas de atención}
\label{sec:met-ablacion}

Se anulan individualmente cada una de las 144 cabezas de atención ($12 \times 12$)
estableciendo su salida a cero (\textit{zero-out}). Para cada cabeza se mide la
caída en F1 macro y en F1 por emoción. Se construye una taxonomía bidimensional
basada en dos métricas ortogonales:

\begin{itemize}[nosep]
    \item Importancia: caída absoluta en F1 macro al ablacionar la cabeza.
    \item Selectividad: concentración del impacto en pocas emociones, medida
    mediante el coeficiente de Gini sobre el vector de caídas por emoción.
\end{itemize}

Las cuatro categorías resultantes son: Critical Specialist (alta importancia, alta
selectividad), Critical Generalist (alta importancia, baja selectividad), Minor
Specialist (baja importancia, alta selectividad) y Dispensable (baja importancia,
baja selectividad).

\subsubsection{Análisis de neuronas FFN}
\label{sec:met-neuronas}

Para cada una de las 36.864 neuronas intermedias del modelo ($12 \times 3.072$) se
extrae la activación post-GELU mediante forward hooks, promediando sobre las posiciones
de la secuencia. Se procesa el conjunto de test completo.

La \emph{selectividad neuronal} para una neurona $n$ y una emoción $e$ se define
como un score tipo Cohen's $d$:
\begin{equation}
    s(n, e) = \frac{\bar{a}_{n,\text{presente}} - \bar{a}_{n,\text{ausente}}}{\sigma_{\text{pooled}}}
    \label{eq:selectividad}
\end{equation}
donde $\bar{a}_{n,\text{presente}}$ es la activación media de la neurona $n$ cuando
la emoción $e$ está presente en el texto, y $\bar{a}_{n,\text{ausente}}$ cuando está
ausente. Se consideran significativas las neuronas con $|s| > 2{,}0$.

Para conectar la especialización neuronal con la compresión SVD, se investiga
una hipótesis concreta: \textit{¿las emociones cuyas neuronas especializadas
caen mayoritariamente en el subespacio eliminado por la SVD son las más
vulnerables a la compresión?} Para responder a esta pregunta, se proyecta el
vector de selectividad de cada emoción (un vector de 3.072 dimensiones por capa,
con la selectividad de cada neurona) sobre los subespacios conservado y eliminado
por la SVD a rango 128, y se calcula la fracción de señal emocional que reside
en cada subespacio. Si la hipótesis es correcta, las emociones con mayor fracción
de señal en el subespacio eliminado deberían mostrar mayor caída de F1 bajo
compresión.

% -------------------------------------------------------------
\subsection{Métricas de evaluación}
\label{sec:met-metricas}
% -------------------------------------------------------------

La evaluación combina métricas a tres escalas: rendimiento global (F1 macro y
micro), impacto diferencial por emoción (F1 por emoción, retención de F1) y
calidad representacional (silhouette score).

\begin{itemize}[nosep]
    \item F1 macro: media no ponderada del F1 por emoción. Trata todas las
    emociones por igual, independientemente de su frecuencia, y es sensible a la
    degradación de emociones minoritarias. Es la métrica principal del proyecto.
    \item F1 micro: F1 calculado sobre el total de predicciones. Favorece
    las emociones frecuentes y proporciona una perspectiva complementaria.
    \item F1 por emoción: permite identificar impactos diferenciales de la
    compresión.
    \item Retención de F1: $F1_{\text{comprimido}} / F1_{\text{baseline}}$,
    expresada como porcentaje.
    \item Ratio de parámetros: $\text{params}_{\text{comprimido}} /
    \text{params}_{\text{original}}$, donde un ratio $< 1$ indica compresión
    efectiva.
    \item Silhouette score: calidad de los clusters en proyecciones t-SNE
    de embeddings \texttt{[CLS]}.
\end{itemize}

\subsection{Entorno de ejecución}
\label{sec:met-entorno}

Todos los experimentos se ejecutaron en Google Colab (GPU Tesla T4, 16~GB VRAM) y/o
en entorno local (CPU). Los notebooks son compatibles con ambos entornos. Se fijó la
semilla aleatoria (42) en PyTorch, NumPy y Python para garantizar la reproducibilidad.
El código se organizó en módulos reutilizables (\texttt{src/data},
\texttt{src/models}, \texttt{src/compression}, \texttt{src/utils}), separados de los
notebooks de experimentación (NB01--NB10).

\subsection{Reproducibilidad}
\label{sec:met-reproducibilidad}

La reproducibilidad del proyecto se aborda en cuatro niveles:

\begin{itemize}[nosep]
    \item \textbf{Código.} El repositorio contiene los nueve notebooks
    de experimentación (\texttt{notebooks/01b}--\texttt{09}) más el
    notebook de \emph{logit lens} (\texttt{notebooks/08b}), todos ejecutables
    end-to-end y con dependencias declaradas en \texttt{requirements.txt}.
    Los módulos compartidos viven en \texttt{src/} (carga de datos,
    SVD, métricas, hooks de extracción de activaciones).

    \item \textbf{Datos.} GoEmotions \cite{demszky2020goemotions} es
    público bajo licencia Apache~2.0 y se descarga vía la librería
    \texttt{datasets} de Hugging Face. Los splits oficiales se utilizan
    sin modificación más allá del filtrado de las cinco emociones
    documentado en \S\ref{sec:met-dataset}.

    \item \textbf{Modelo.} El checkpoint baseline
    (\texttt{bert-goemotions-23emo-final}, 110~M parámetros) se obtiene
    aplicando el procedimiento de \emph{fine-tuning} descrito en
    \S\ref{sec:met-modelo} sobre BERT-base-uncased
    \cite{devlin2019bert} con la semilla fija. La ejecución completa
    ocupa unas 25~minutos en una GPU T4 de Colab.

    \item \textbf{Resultados intermedios.} Cada notebook exporta sus
    resultados a \texttt{csvs/notebookN/} y \texttt{results/notebookN/}
    como ficheros CSV planos (más algunos PNG). Esto permite reproducir
    cualquier figura o tabla de la memoria sin tener que re-ejecutar
    los experimentos: basta con cargar los CSV correspondientes. Los
    análisis estadísticos posteriores (correlaciones de Spearman,
    bootstrap por emoción, intervalos de confianza al 95\%) son
    deterministas dada la semilla y operan sobre estos CSV.
\end{itemize}

La principal fuente de no-determinismo residual es el orden de los
\textit{batches} dentro de cada época de \emph{fine-tuning}, que en GPU
introduce variaciones en la cuarta cifra decimal del F1 macro entre
ejecuciones independientes. Las cifras reportadas en la memoria
corresponden a la ejecución con semilla 42, y cualquier divergencia
respecto a una réplica posterior debería caer dentro de ese margen.